% Method section. Written from the code on the campaign branch (a936cd63a), not from older docs.
% Every behavioural claim carries a "% src: path:line" comment. Paths are relative to the repo root; the
% agent-optimize scripts live in skills/algorithms/agent-optimize/scripts (abbreviated AO/ below).
% No result numbers: \TODO{} placeholders only.
% Requires: \usepackage[T1]{fontenc}, amsmath, amssymb, booktabs, algorithm, algorithmic, natbib. \TODO is provided here if the main file lacks it.
\providecommand{\TODO}[1]{\textbf{[TODO: #1]}}
\newcommand{\artifact}{a}
\newcommand{\Dopt}{\mathcal{D}_{\mathrm{opt}}}
\newcommand{\Dsel}{\mathcal{D}_{\mathrm{sel}}}
\newcommand{\Dfin}{\mathcal{D}_{\mathrm{fin}}}
\newcommand{\rbar}{\bar r}
\newcommand{\SE}{\mathrm{SE}}

\section{Method}
\label{sec:method}

\subsection{Problem formalization}
\label{sec:method:problem}

\emph{Frozen model, editable artifact.}
The target model weights $\theta$ are frozen throughout. What changes is an \emph{artifact} $\artifact\in\mathcal{A}$, a
set of files from which the agent harness is assembled at run time: the system prompt, agent-side tool definitions
(code and schemas), MCP server tool descriptions, and skill packages (instruction files plus helper scripts).
% src: core/cap_evolve/harness.py:318 (evaluate_candidate runs the adapter on a candidate directory); the framework
% treats the candidate as an opaque directory of files that the proposer edits and the adapter loads.
A search method sees $\theta$ only through rollouts: running the harness built from $\artifact$ on a task $x$ with
seed $\omega$ yields a scalar reward $r(\artifact,x,\omega)\in[0,1]$ from the benchmark's scorer, and a resource
vector (tokens, dollars, wall time).
% src: core/cap_evolve/harness.py:318-640 (per-trial rollouts, scores and costs persisted per task and trial)

\emph{Tasks and splits.}
Let $\mathcal{D}$ be a distribution over tasks. We use three disjoint task sets, drawn once and frozen by a split seed:
an optimization set $\Dopt$ whose rollouts the proposer may read, a selection set $\Dsel$ on which candidates are
compared and promoted, and a final set $\Dfin$ that is scored once per run at the end.
% src: core/cap_evolve/splits.py:26-60 (frozen splits.json with train/val/test and a test_used seal)
The shipped hill-climb and agent-optimize implementations use the validation split for \emph{both} roles: the
proposer reads validation rollouts and per-task rewards, and the same split gates acceptance, so
$\Dopt=\Dsel=\mathrm{val}$ and the train split is unused. This adaptive reuse is a property of the implementation that
we disclose rather than hide; the competing methods in our comparison use the train split (80 tasks) for feedback and the
validation split (40 tasks) for selection, so the stock CapEvolve arms differ in data exposure, which we report as such; its consequences are discussed in Section~\ref{sec:method:threat}.
% src (competitors): research/capevolve-paper/decisions.md (stock CapEvolve uses VAL traces for feedback and VAL for gating; competitors use train-80 feedback + val-40 selection)
% src: core/cap_evolve/harness.py:2512 (parent's val trajectories copied into workdir/trajectories);
% core/cap_evolve/harness.py:4591-4592 (instructions built from current_val per-task rows);
% research/capevolve-paper/audit/defects.csv (D03: val is both feedback source and acceptance split)
Write $\rbar_n(\artifact,x)=\frac1n\sum_{j=1}^{n}r(\artifact,x,\omega_j)$ for the task-level mean over $n$ trials and
$J_{S}(\artifact)=\frac1{|S|}\sum_{x\in S}\mathbb{E}_\omega\, r(\artifact,x,\omega)$ for the mean reward on a set $S$.
The quantity of scientific interest is $J_{\Dfin}$ of the artifact the search returns.

\emph{Objective and budget.}
A search receives a budget $B$ that caps evaluation calls, dollars, optimizer dollars, iterations, and consecutive
rejects; the loop stops when any cap is reached.
% src: core/cap_evolve/rundir.py:368-383 (budget_exhausted: stop_at_reward, max_iterations, max_metric_calls, max_usd,
% max_optimizer_usd, stall)
The goal is
\begin{equation}
  \max_{\artifact\in\mathcal{A}}\; J_{\Dfin}(\artifact)
  \quad\text{subject to}\quad C_{\mathrm{opt}}(\text{search})\le B,
\end{equation}
where $C_{\mathrm{opt}}$ counts \emph{every} evaluation the search spends (screening, confirmation, merge verification,
probes) plus proposer cost. The method returns the last promoted candidate, never the best-looking candidate on
$\Dsel$. Deployment cost per task is reported separately and is not part of $B$.
% src: core/cap_evolve/harness.py:3315-3316 (set_best only on accept); core/cap_evolve/harness.py:4676-4752 (finalize
% scores best_dir, not a re-selected maximum)

\subsection{Two search algorithms in one notation}
\label{sec:method:algs}

Both algorithms follow the propose--evaluate--promote pattern of reflective prompt and skill optimizers \citep{gepa,skillopt}. They maintain a set of candidate artifacts with parent links and an \emph{incumbent} $\artifact^\star$
together with the incumbent's measurement $m^\star$ (per-task means on $\Dsel$). They differ in who proposes, what the
proposer sees, how parents are chosen, and how evaluation effort is allocated; they can share one promotion rule
(Section~\ref{sec:method:P}). Table~\ref{tab:method:algs} (Appendix~\ref{app:skeleton}) summarizes the differences and
Algorithm~\ref{alg:skeleton} (Appendix~\ref{app:skeleton}) gives the shared skeleton.


% src (hill-climb): core/cap_evolve/harness.py:4581 (budget check), 4591 (parent = run_dir.best_id), 4592 (instructions),
%   4597 (convergence monitor), 4608 (run_step), 4618 (incumbent update); core/cap_evolve/gate.py:422-459 (paired gate, bar = k_se*SE);
%   core/cap_evolve/harness.py:3199-3207 (footprint then _paired_deltas); research/capevolve-paper/audit/algorithm_audit.md:94-102.
% src (agent-optimize): AO/act.py:183-200 (propose, parent from --parent or champion, dag_parallel forces champion),
%   AO/act.py:329 (probe), AO/act.py:388-411 (judge: reward_gated when cost_gating on, posterior otherwise),
%   AO/act.py:512 (promote), AO/act.py:556 (prune), AO/act.py:583 (merge), AO/act.py:749 (finalize), AO/act.py:847 (VERBS);
%   AO/SKILL.md:15-21 ("never view test traces" is an instruction); core/cap_evolve/posterior.py:1-10.
% Default k_se: the CLI forwards spec gate_k_se (default 1.0): research/capevolve-paper/audit/algorithm_audit.md:15-20.


% src: core/cap_evolve/harness.py:4581-4618 (loop), 3093-3096 and 3130 (integrity snapshot and verify around the proposer),
%   core/cap_evolve/harness.py:3118-3150 (a tampered step is discarded through _tamper_step, not scored 0),
%   core/cap_evolve/harness.py:4676-4752 (finalize scores best and seed on test, then commit_test).

\emph{Hill-climb.}
Each iteration forks the incumbent into a working directory, injects its validation trajectories and a focus
description (all failing tasks, or one task in a cyclic or hardest-first schedule), runs the proposer subprocess, and
evaluates the result on the full validation split with $n$ trials per task. The per-task rewards of candidate and
incumbent give a paired delta vector; the step is accepted by the gate described below and, on accept, the candidate
becomes the next parent.
% src: core/cap_evolve/harness.py:4591-4618; core/cap_evolve/harness.py:3158-3160 (evaluate_candidate on val);
%   core/cap_evolve/harness.py:3207 (_paired_deltas); core/cap_evolve/harness.py:3315-3316 (set_best)
Unless disabled, the delta vector is first restricted to the candidate's \emph{footprint}: tasks whose persisted
rollouts mention none of the surfaces the edit touched enter as exact zeros. The mask is dropped (the vector keeps its
length) whenever the footprint cannot be established, and it is the only component of the default gate that uses
information from rollout text.
% src: core/cap_evolve/footprint.py:1-80 (abstains/returns None when it cannot localize); core/cap_evolve/harness.py:3199-3207
The default decision is the paired one-SE test: with $\Delta_t=\rbar'_t-\rbar^\star_t$ over tasks scored on both sides,
accept iff $\bar\Delta > k_{\mathrm{se}}\,\SE(\Delta)$.
% src: core/cap_evolve/gate.py:422-459 (mean_d > k_se*se; SE = sqrt(var/n) across tasks)

\emph{Agent-optimize.}
A driver agent owns the search; scripts decide and record. After a digest, the driver proposes an edit for a
failure-cluster hypothesis, creating a node whose parent may be any existing node. Evaluation is lazy: a \emph{probe}
runs only the (task, trial) cells missing from an evidence ledger keyed by the content hash of the candidate, so
byte-identical candidates share evidence. \texttt{promote} requires full-split ledger coverage and then calls a judge.
% src: AO/act.py:183-246 (propose, 2-call protocol); AO/act.py:329-385 (probe over missing ledger cells);
%   core/cap_evolve/eval_index.py:55-80 (counts/missing pooled by cap_hash); AO/act.py:512-535 (coverage check then judge)
Without opt-in switches the judge is a Bonferroni-corrected $z$-test on the reward difference with a fixed SE floor
(\texttt{reward\_gated}), and a Beta-posterior rule is used only when cost gating is switched off; with
\texttt{promotion\_\allowbreak rule: P} the judge is the same function hill-climb calls.
% src: AO/act.py:400-411; core/cap_evolve/objectives.py:41,84,166-176 (NOISE_FLOOR=0.031 floors the SE);
%   core/cap_evolve/posterior.py:17-19 (DMIN, P_ACC, canary constants); AO/act.py:512-520 (judge_p when promotion_params)
\texttt{finalize} seals and scores the test split once.
% src: AO/act.py:749-845

\subsection{Promotion rule P}
\label{sec:method:P}

Because the two default gates have different effective bars, comparing the algorithms would compare gate strictness as
well as search. Rule P is a single function both algorithms can call (spec key \texttt{promotion\_\allowbreak rule: P}, default
off; implemented as the shared module of issues \#785/\#786). It has two stages and four parameters
$(k,\ n_{\mathrm{conf}},\ \delta_{\min},\ z_{\mathrm{conf}})$ plus a coverage floor $\rho$.
% src: core/cap_evolve/promotion.py:28-52 (PromotionParams: k_se=1.0, n_conf=10, dmin=0.02, z_conf=2.0, min_coverage=0.6)

\begin{algorithm}[t]
\caption{Promotion rule P for candidate $\artifact'$ against incumbent $\artifact^\star$.}
\label{alg:P}
\begin{algorithmic}[1]
\REQUIRE screen measurements $m',m^\star$ on $\Dsel$; $k,\ n_{\mathrm{conf}},\ \delta_{\min},\ z_{\mathrm{conf}},\ \rho$
\STATE $\Delta_t\gets m'_t-m^\star_t$ for $t$ scored on both sides (\emph{unmasked}: no footprint restriction)
\IF{$|\Delta|<2$ \textbf{or} $\bar\Delta\le 0$ \textbf{or} $\bar\Delta\le k\,\SE(\Delta)$} \STATE \textbf{return} reject \ENDIF
\STATE draw one trial offset $o$ (agent-optimize: beyond every trial either artifact has; hill-climb: $o=1000+i\,n_{\mathrm{conf}}$ at iteration $i$); fix seeds as $\mathrm{base}+o+j$
\STATE $m'_{\mathrm{c}}\gets \mathrm{Eval}(\artifact',\Dsel,n_{\mathrm{conf}};o)$;\quad $m^\star_{\mathrm{c}}\gets \mathrm{Eval}(\artifact^\star,\Dsel,n_{\mathrm{conf}};o)$
\STATE $\Delta^{c}_t\gets m'_{\mathrm{c},t}-m^\star_{\mathrm{c},t}$;\quad $\mathrm{cov}\gets |\Delta^{c}|/|\Delta|$
\IF{$\mathrm{cov}<\rho$ \textbf{or} $\bar\Delta^{c}<\delta_{\min}$ \textbf{or} $\bar\Delta^{c}/\SE(\Delta^{c})<z_{\mathrm{conf}}$} \STATE \textbf{return} reject \ENDIF
\RETURN promote, with new baseline $m^\star\gets m'_{\mathrm{c}}$
\end{algorithmic}
\end{algorithm}
% src: core/cap_evolve/promotion.py:73-92 (screen, confirm), 111-130 (decide: confirm_fn is called only after a screen pass);
%   core/cap_evolve/harness.py:3064-3069 (P forces footprint off and paired mode), 3270-3299 (fresh tags <cid>__conf and
%   <cid>__confbase, trial_offset=fresh_from, fixed_seeds=True, rejection recorded in decision.reason), 4618 (next baseline
%   = baseline_val); hill-climb offset fresh_from = 1000 + iterations*n_conf at harness.py:3273; agent-optimize offset = _first_free_trial
%   (AO/act.py:443-450); core/cap_evolve/harness.py:371-376 (fixed_seeds: seed is exactly base+offset+k, ignoring the ledger);
%   AO/act.py:463-510 (judge_p: same promotion.decide, common offset from _first_free_trial, budget checked before each
%   fresh evaluation, champion's first fresh trial index stored in promotion_p.json for the next screen).

When the paired SE is exactly zero (every task moved identically) and the mean is positive, $z$ is taken as infinite and the stage passes, subject to $\delta_{\min}$ and coverage; a mean of zero or less never passes, and fewer than two paired tasks never passes.
% src: core/cap_evolve/promotion.py:67-79 (_z: +inf when se==0 and mean>0; screen requires n>=2 and mean>0), 80-92 (confirm)
Two design points matter for the measurement contract. First, the confirmation is computed from trials that did not
participate in the screen, so the reported confirmation delta is not inflated by selecting the candidate on the same
draws. Second, the \emph{confirmation} of the promoted candidate, not its screen measurement, is the next baseline,
which removes the selection bias from the reference that later candidates must beat; on resume the confirmation
measurement is read back.
% src: core/cap_evolve/harness.py:1128-1135 (resume_val reads <tag>__conf when rule P is on); core/cap_evolve/harness.py:4618
Parameter values used in the experiments: $k=1$, $n_{\mathrm{conf}}=3$, $\delta_{\min}=0$,
$z_{\mathrm{conf}}=1.282$ (the one-sided $90\%$ normal quantile), $\rho=0.95$ (the library defaults are $n_{\mathrm{conf}}=10$,
$\delta_{\min}=0.02$, $z_{\mathrm{conf}}=2$, $\rho=0.6$). We chose them by a budget-exact simulation of one search per suite, using the difficulty observed on the seed artifact (Appendix~\ref{app:sim}). The error rate is per search, not per edit: in an all-null search of about ten candidates, rule P falsely promotes at least once with probability \simPFalseA{}, \simPFalseB{} and \simPFalseC{} on Suites A, B and C, about one percent per candidate. The one-stage $k=1$ gate does so with probability \simKoneFalseA{}, \simKoneFalseB{} and \simKoneFalseC{}. Rule P promotes a single candidate with a true gain of $+0.08$ with probability \simPPowerEightA{}, \simPPowerEightB{} and \simPPowerEightC{} on Suites A, B and C, against \simKonePowerEightA{}, \simKonePowerEightB{} and \simKonePowerEightC{} for the one-stage gate. Rule P gives up power by design. Appendix~\ref{app:sim} explains why we did not choose five confirmation trials.
% src (numbers): research/capevolve-paper/analysis/out/sim_budget.md (rows "f: one-stage k=1", "c: P n_conf=3 z=1.282 dmin=0",
%   "a: P n_conf=3 z=2.0 dmin=.02"; power table), produced by research/capevolve-paper/analysis/scripts/sim_budget.py; these are
%   simulation outputs under that script's generative assumptions, not measurements on the benchmark.
% src (library defaults): core/cap_evolve/promotion.py:28-35
Both \capevolve{} arms judge through rule P with the footprint mask off and joint coverage, so the two algorithms share one
acceptance rule. The confirmation costs $2\,|\Dsel|\,n_{\mathrm{conf}}$ rollouts per candidate that passes the screen and
nothing for candidates that do not; these rollouts are charged to the budget $B$.
% src (values): frozen E10 configuration communicated by the campaign lead; research/capevolve-paper/decisions.md (gate parity entry)
% src: core/cap_evolve/harness.py:3270-3300; test core/tests/test_promotion_rule_p.py::test_confirmation_cost_is_charged_to_the_budget
% (spent.metric_calls grows by exactly 2*T*n_conf); research/capevolve-paper/decisions.md (campaign parameter choice)

\emph{Experimental configuration (frozen before the search runs).}
Screening evaluates each validation task with one trial on the 40 validation tasks. Each search is limited to
$600$ target rollouts, counting every evaluation including confirmations, and at most $4\mathrm{M}$ non-cache-read
optimizer tokens. The target is Gemma-4-31B-it with a limit of $30$ turns and temperature $0$. The reward is the
LibreOffice-25.8 recalculated hard score. agent-optimize runs its current digest-and-act engine with its stock switches (parallel branching, smart merge,
cost gating, failure clustering, pregate and context digest on), and adds active evaluation and fresh-trial-only confirmation.
% src (values): frozen E10 configuration communicated by the campaign lead; research/capevolve-paper/decisions.md
%   ("Frozen E10 search parameters"); switch defaults: core/cap_evolve/optimizer_config.py DEFAULTS (dag_parallel True, active_eval False,
%   smart_merge True, cost_gating True, failure_clustering True, pregate True, context_digest True)

\subsection{Measurement contract}
\label{sec:method:contract}

The statistics above are only as good as the measurements they consume. We state the contract each run must satisfy and,
for each clause, where it is enforced and where it is only a deployment requirement.

\emph{Evaluator isolation.}
(i) \emph{Gold-free sandbox.} The agent's sandbox must not contain the gold answers. This is benchmark-specific and
lives in the adapter: for SpreadsheetBench the default mounts a gold-free copy of the dataset (inputs and a writable
output directory), reads gold from the pristine directory at scoring time, and refuses a mount that contains
gold-named, unrecognized or symlinked files; mounting gold requires two explicit opt-ins.
% src: templates/adapters/spreadsheetbench/adapter.py:138-148 (AGENT_VIEW default on), 285-348 (allowlist, symlink refusal),
%   421-444 (refusal messages); core/tests/test_spreadsheetbench_gold_isolation.py (score equivalence with and without gold in the mount)
(ii) \emph{Proposer read isolation.} The proposer is a coding-agent subprocess that inherits the environment and has
shell access; the library does not sandbox it. Read isolation of gold, split files and test results from the proposer is
therefore a deployment requirement, which our experiments satisfy by running the proposer under a separate user id and
verifying denial with a negative test before launch; this is a property of the harness around the library, not of the
algorithms (Appendix~\ref{app:incidents}).
% src: research/capevolve-paper/audit/algorithm_audit.md:55 and :182 (unsandboxed subprocess inheriting os.environ);
%   research/capevolve-paper/decisions.md (INTEGRITY entry: separate uid, negative test).
(iii) \emph{Tamper detection.} Optionally, a content-hash manifest over the eval surface (gold, task data, tests,
scorers, adapters, split files) is taken before the proposer runs, the files are set read-only, and the manifest is
verified before the gate; a changed, added or deleted protected file voids the measurement. A voided candidate is
discarded, not scored $0$, because a tampered number describes a compromised harness rather than the edit.
% src: core/cap_evolve/integrity.py:1-60 (semantics), 65-77 (DEFAULT_PATTERNS); core/cap_evolve/harness.py:3093-3096, 3130-3150
Detection is host-independent but it is detection after the fact and does not stop reads.

Three further clauses (errored-trial policy, joint coverage, fresh-trial confirmation) are specified in Appendix~\ref{app:contract}; the per-arm settings are in Appendix~\ref{app:config}.

\emph{Test seal and the campaign-level ledger.}
Within a run, the final split is sealed: finalize reserves the seal, refuses a second attempt once test rollouts exist,
scores the returned artifact (the last promoted candidate) and the seed on the same test tasks, and burns the seal only after the results are written.
% src: core/cap_evolve/rundir.py:428-477 (reserve_test, begin_test_attempt, commit_test); core/cap_evolve/harness.py:4707-4752
The seal is \emph{per run directory}. A campaign that repeats runs on one final split would therefore consume it
repeatedly, and any decision informed by an earlier final score is adaptive reuse of the holdout
\citep{adaptive-data-analysis-holdout-reuse}. We hence keep a campaign-level ledger keyed by the hash of the sorted
final-split task ids, recording run, arm and time, and we report the reuse count with every headline number.
% src: research/capevolve-paper/audit/defects.csv (D07: test_used is per run; reuse_baseline resets it; no cross-run ledger);
%   research/capevolve-paper/audit/issue_drafts/D07.md:21 (proposal). The ledger is a campaign-protocol record, not library code
%   at this commit. The exposure table is research/capevolve-paper/audit/exposure_ledger.csv.

\subsection{Threat model and what is not guaranteed}
\label{sec:method:threat}

We make no statistical guarantee for the search-time rules; they are heuristics with measured operating characteristics.
The paired one-SE gate accepts when $\bar\Delta>\SE$; under a normal approximation and a true null this fires with
probability $\Pr(Z>1)\approx0.16$ per test, and repeated looks compound it (in the budget-exact simulation of Section~\ref{sec:method:P}, an all-null search of about ten candidates falsely promotes at least once with probability \simKoneFalseB{} to \simKoneFalseA{} under the one-stage $k=1$ gate, depending on the suite).
% src: core/cap_evolve/gate.py:451-452; the 0.16 is the standard normal tail Pr(Z>1)=0.1587. Simulated rates:
%   research/capevolve-paper/analysis/out/sim_promotion.md, sim_gate.md
The Beta-posterior rule is a Bayesian decision rule under a modelling assumption; we make no frequentist claim for it,
and its sequential-validity wording in the code comments is not an established bound.
% src: core/cap_evolve/posterior.py:1-10; research/capevolve-paper/audit/defects.csv (D10)
Rule P lowers the false-promotion rate but does not make it zero: the confirmation controls trial noise (seed-level luck),
including the selection of a candidate on its own screen draws, and the unmasked vector removes the footprint mask as a source of
inflated deltas. It does not control same-task luck: an edit that is genuinely better on \emph{these} tasks passes both stages. It does \emph{not} control adaptivity over the task set: the proposer reads validation traces and the
same split gates acceptance, so a candidate can overfit the validation tasks and pass both stages. Only the sealed final
split gives a number that is not selected on, and only to the degree that it has not been reused
\citep{adaptive-data-analysis-holdout-reuse}. Anytime-valid sequential designs
\citep{confidence-sequences,shawn2026pace} would address repeated looks and are not implemented in the search-time gates.

The proposer is an LLM agent with shell access to the file system the evaluation uses. The system provides tamper detection around each proposal (when enabled), a gold-free agent sandbox in the SpreadsheetBench adapter, a per-run test seal, and logs of both the gate statistic and its unmasked version. It does not sandbox the proposer (a deployment requirement), stop it from reading test outputs left on disk by earlier runs, detect adapter feedback that leaks gold, or guarantee that a promoted edit transfers beyond the final tasks.
% src: research/capevolve-paper/audit/gold_access_scan.md; core/cap_evolve/adapter.py:104; research/capevolve-paper/audit/algorithm_audit.md:182
Finally, the benchmarks we use were partly exposed to earlier development; which final sets are development-exposed is
tabulated in \texttt{audit/exposure\_ledger.csv} (Section~\ref{sec:problem}) and the headline claims are restricted to sets that were not.
% src: research/capevolve-paper/decisions.md (development-exposed sets); research/capevolve-paper/audit/exposure_ledger.csv
